\documentclass[../../../main.tex]{subfiles}

\begin{document}

\section{Recurrence Relation}

Another very popular type of problems in math olympiads and
competition in general is recurrence relation. This section discusses
different forms of recurrence relation and related problems. One more
thing before we get started is that this section is largely based on
a series of YouTube videos from WinnersKMO \cite{MATH01-5}.

As you probably know, recurrence relation focuses on $n^\text{th}$
term given its relation with preceding $k^\text{th}$ terms.
One very famous example would be Fibonacci sequence whose terms
are defined as $F_n = F_{n-1} + F_{n-2}$ where $F_0 = F_1 = 1$.
The problems are relatively easy in my opinion until it gets
incredibly hard, so let's get started with some easy cases.

{\Large\textcolor{BrickRed}
{\textbf{I. $\mathbf{a_n = a_{n-1} + f(n)}$}}}

This default type has an effective way of solving because it always
gets down to a single formula, which becomes quite evident when we
telescope. Consider the following equations.
\begin{align*}
    a_n - a_{n-1} &= f(n) \\
    a_{n-1} - a_{n-2} &= f(n-1) \\
    &\vdots \\
    a_2 - a_1 &= f(2)
\end{align*}
Adding the equations, the left-hand side telescopes and we obtain
$a_n - a_1 = \sum_{k=2}^n f(k)$.

\begin{important}
Given a recurrence relation $a_n = a_{n-1} + f(n)$, each term can be
represented with the following formula.
\[
a_n = \sum_{k=1}^n f(k) + a_1 - f(1)
\]
\end{important}

There will be lots of formulas from now, but I honestly think there
is no need to memorize them because if we just remember the idea to
get to the formulas, then we can solve the problems. Here is an
example problem before we move on to the next case.

\begin{problem}{Enoch's Problem}{}
Let $a_1, \ldots, a_k$ be a sequence such that $a_n = a_{n-1} +
n^2 - 2n + 1$ for $n > 1$ and $a_1 = 2$. Find $a_{101}$.
\end{problem}
\begin{solution}
This sequence has the form $a_n = a_{n-1} + f(n)$ where $f(n) = n^2
- 2n + 1$. Telescoping, we obtain the following equations.
\begin{align*}
    a_{101}
    &= \sum_{k=1}^{101} (k-1)^2 + 2 - 1
    = \sum_{n=0}^{100} n^2 + 1 \\
    &= \frac{100(100 + 1)(200 + 1)}{6} + 1
    = 338351
\end{align*}
Therefore, $a_{101} = 338351$.
\end{solution}

Continuing with the first form, here is the second one.

{\Large\textcolor{BrickRed}
{\textbf{II. $\mathbf{a_n = a_{n-1} f(n)}$}}}

This case is also very popular and can be solved using the same method
as before.
\begin{align*}
    \frac{a_n}{a_{n-1}} &= f(n) \\
    \frac{a_{n-1}}{a_{n-2}} &= f(n-1) \\
    &\vdots \\
    \frac{a_2}{a_1} &= f(2)
\end{align*}
Telescoping, we obtain $\frac{a_n}{a_1} = \prod_{k=2}^n f(k)$.

\begin{important}
The terms in a recurrence relation $a_n = a_{n-1} f(n)$ can be
represented with the following formula.
\[
a_n = a_1 \prod_{k=2}^n f(k)
\]
\end{important}

Here is another problem derived from the previous one.

\begin{problem}{Enoch's Problem}{}
Let $a_1, \ldots, a_k$ be a sequence such that $a_n = a_{n-1} (
n^2 - 2n + 1)$ for $n > 1$ and $a_1 = 1$. Find $a_{101}$.
\end{problem}
\begin{solution}
Notice that this sequence has the form $a_n = a_{n-1} f(n)$ where
$f(n) = (n-1)^2$. Telescoping, we obtain the following equations.
\[
a_{101}
= \prod_{k=2}^{101} (k-1)^2
= \prod_{n=1}^{100} n^2
= (100!)^2
\]
Therefore, $a_{101} = (100!)^2$.
\end{solution}

The first two cases were relatively simple. Continuing, we can
continue with their variations.

{\Large\textcolor{BrickRed}
{\textbf{III. $\mathbf{a_n = p a_{n-1} + q}$}}}

From this case, there are many different ways to get an answer.
Let's start with my personal favorite one.

{\large\textbf{Solution I.}}

The first method is to directly substitute.
\begin{align*}
    a_n
    &= p a_{n-1} + q \\
    &= p (p a_{n-2} + q) + q
    = p^2 a_{n-2} + (p + 1)q \\
    &= p^2 (p a_{n-3} + q) + (p + 1)q
    = p^3 a_{n-3} + (p^2 + p + 1)q \\
    &\qquad\quad\vdots \\
    &= p^{n-1} a_1 + q \sum_{k=0}^{n-2} p^k
\end{align*}
Rearranging the equation, we obtain the following equation.
\[
a_n = p^{n-1} a_1 + \frac{1 - p^{n-1}}{1 - p} \cdot q
\]
I personally like this method because it feels intuitive and quite
natural to just substitute. The second method is slightly less
intuitive, but it should be faster and better if you prone to make
arithmetic mistakes.

{\large\textbf{Solution II.}}

This method will convert the form of this recursive equation.
First, assume that there exists an $\alpha$ such that the following
equation holds.
\[
a_n - \alpha = p (a_{n-1} - \alpha)
\]
If there exists such alpha, then it is evident that we can telescope
to obtain the following equation.
\[
a_n - \alpha = p^{n-1} (a_1 - \alpha)
\]
This is much easier to work with compared to our initial equation.

From the equation $a_n - \alpha = p (a_{n-1} - \alpha)$, we can notice
that $p a_{n-1} + q - \alpha = p (a_{n-1} - \alpha)$ and
$q = \alpha (1 - p)$. In other words, $\alpha = \frac{q}{1 - p}$.
This shows that $\alpha$ always exists given that $p \neq 1$.
In other words, whenever we see this case, we can let $\alpha =
\frac{q}{1 - p}$ and use $a_n - \alpha = p^{n-1} (a_1 - \alpha)$.
Substituting,
\begin{align*}
    a_n
    &= p^{n-1} a_1 + \alpha \left( 1 - p^{n-1} \right)
    = p^{n-1} a_1 + \frac{q}{1 - p} \left( 1 - p^{n-1} \right) \\
    &= p^{n-1} a_1 + q \sum_{k=0}^{n-2} p^k
    = p^{n-1} a_1 + \frac{1 - p^{n-1}}{1 - p} \cdot q
\end{align*}
is obtained and this result is consistent.

Before moving on to the next case, here is an example problem.

\begin{problem}{Enoch's Problem}{}
Find $a_{101}$ if $a_n = 2a_{n-1} + 3$ for $n > 1$ and $a_1 = 2$.
\end{problem}
\begin{solution}
First, we can try substituting $a_{n-1}$ in the given equation in
terms of $a_{n-2}$ and continue the process until we reach $a_1$.
\begin{align*}
    a_{101}
    &= 2 a_{100} + 3 \\
    &= 2 (2 a_{99} + 3) + 3
    = 2^2 a_{99} + (2 + 1) 3 \\
    &\qquad\quad\vdots \\
    &= 2^{100} a_1 + 3 \sum_{k=0}^{99} 2^k
\end{align*}
Therefore, $a_{101} = 2^{101} + \frac{1 - 2^{100}}{1 - 2} \cdot 3
= 2^{101} - 3 + 3 \cdot 2^{100} = 2^{100} 5 - 3$.
\end{solution}

Continuing from this case, we have its extension where $q$ is not
a constant but a function.

{\Large\textcolor{BrickRed}
{\textbf{IV. $\mathbf{a_n = k a_{n-1} + f(n)}$}}}

This problem too can be solved with very similar methods from the
previous, its just a little more unkind. Let's start with my
favorite method.

{\large\textbf{Solution I.}}

In this solution, we will define a new term $b_n = \frac{a_n}{k^n}$.
Consider the following equations.
\begin{align*}
    \frac{a_n}{k^n} &= \frac{a_{n-1}}{k^{n-1}} + \frac{f(n)}{k^n} \\
    b_n
    &= b_{n-1} + \frac{f(n)}{k^n}
    = b_{n-1} + g(n)
\end{align*}
This is the same as Case I and telescoping gives us
\[
b_n = \sum_{k=1}^n g(k) + b_1 - g(1)
\]
and rewriting this, we obtain the following.
\[
\frac{a_n}{k^n}
= \sum_{l=1}^n \frac{f(l)}{k^l} + \frac{a_1}{k} - \frac{f(1)}{k}
= \sum_{l=1}^n \frac{f(l)}{k^l} + \frac{a_1 - f(1)}{k}
\]
This method simply defines a new sequence to represent it into an
easier problem.

{\large\textbf{Solution II.}}

The second method is very similar to the second solution for Case III.
First, we will assume that there exists a function $F(n)$ such that
the following equation holds.
\[
a_n - F(n) = k(a_{n-1} - F(n-1))
\]
Rewriting this, we can see that $a_n = k^{n-1} (a_1 - F(1)) + F(n)$.
From here it depends on the problem, but it shouldn't be too difficult
because we only need to find $F(n)$.

Before moving on to the next case, here is a practice problem.

\begin{problem}{Enoch's Problem}{}
Find $a_n$ in terms of $n$ if $a_n = 3a_{n-1} - 2n + 1$ and $a_1 = 2$.
\end{problem}
(\href{https://youtu.be/meqmlhLmCFk}{Video Solution})
\begin{solution}
First, we can try to divide both sides by $3^n$ and let
$b_n = \frac{a_n}{3^n}$.
\begin{align*}
    \frac{a_n}{3^n} &= \frac{a_{n-1}}{3^{n-1}} - \frac{2n - 1}{3^n} \\
    b_n &= b_{n-1} - \frac{2n - 1}{3^n}
\end{align*}
Telescoping we obtain $b_n - b_1 = -\sum_{k=2}^n \frac{2k - 1}{3^k}$.
Let $S = \sum_{k=2}^n \frac{2k - 1}{3^k}$ and consider the following
equations for index-shifting.
\begin{align*}
    S &= \sum_{k=2}^n \frac{2k - 1}{3^k} \\
    \frac{S}{3}
    &= \sum_{k=2}^n \frac{2k - 1}{3^{k+1}}
    = \sum_{j=3}^{n+1} \frac{2j - 3}{3^j} \\
    \frac{2}{3} S
    &= \sum_{k=2}^n \frac{2k - 1}{3^k} - \left(
        \sum_{j=2}^n \frac{2j - 3}{3^j}
            - \frac{1}{9} + \frac{2n - 1}{3^{n+1}}
    \right) \\
    &= \sum_{k=2}^n \frac{2}{3^k} + \frac{1}{9}
        - \frac{2n - 1}{3^{n+1}}
\end{align*}
Therefore, $S = \sum_{k=2}^n \frac{1}{3^{k-1}} + \frac{1}{6} -
\frac{2n - 1}{2 \cdot 3^n} =
\frac{1}{3} \cdot \frac{ 1 - \frac{1}{3^{n-1}} }{ 1 - \frac{1}{3} } +
\frac{1}{6} - \frac{2n - 1}{2 \cdot 3^n}$. Substituting,
\begin{align*}
    \frac{a_n}{3^n} - \frac{2}{3} &= -\left(
        \frac{1}{2} - \frac{1}{2 \cdot 3^{n-1}}
        + \frac{1}{6} - \frac{2n - 1}{2 \cdot 3^n}
    \right) \\
    \frac{a_n}{3^n}
    &= -\frac{1}{2} + \frac{1}{2 \cdot 3^{n-1}}
        - \frac{1}{6} + \frac{2n - 1}{2 \cdot 3^n} + \frac{2}{3} \\
    &= \frac{1}{2 \cdot 3^{n-1}} + \frac{2n - 1}{2 \cdot 3^n}
    = \frac{1}{2 \cdot 3^{n-1}} + \frac{n}{3^n}
        - \frac{1}{2 \cdot 3^n}
\end{align*}
Therefore, $a_n = n + 1$.
\end{solution}

Continuing from this case, here is the next one. From this case
the problems could get more complicated, but I don't think they are
necessarily harder because they build from the four cases that we
discussed so far.

{\Large\textcolor{BrickRed}
{\textbf{V. $\mathbf{a_n = \frac{p a_{n-1}}{q a_{n-1} + r}}$}}}

This type of problem is relatively simple. To solve them,
we can first take the reciprocal of both sides to obtain a more
simple problem.
\[
\frac{1}{a_n}
= \frac{q a_{n-1} + r}{p a_{n-1}}
= \frac{q}{p} + \frac{r}{p a_{n-1}}
\]
Now substituting $b_n = \frac{1}{a_n}$, we obtain Case III!
The next case is a little more complex, but is also quite common.

{\Large\textcolor{BrickRed}
{\textbf{VI. $\mathbf{p a_{n+1} + q a_n + r a_{n-1} = 0}$}}}

For this case, we can separate the cases into two cases. If the
problem maker was kind, we could sometimes get $p + q + r = 0$.
In that case, we can simply substitute $q = -p - r$ to obtain
the following equation.
\[
p (a_{n+1} - a_n) = r (a_n - a_{n-1})
\]
Letting $b_n = a_n - a_{n-1}$, it is now trivial.

If the problem is not nice enough, which is the most likely case,
we will get $p + q + r \neq 0$. This now becomes quite complicated.
For this case, we can try dividing both sides by $p$. Notice that
because $\frac{q}{p}$ and $\frac{r}{p}$ are constants, we can always
find $\alpha$ and $\beta$ such that the following equation holds
with Vieta's formula.
\[
a_{n+1} - (\alpha + \beta) a_n + \alpha\beta a_{n-1} = 0
\]
From here, we can again split the equation into two cases depending
on the values of $\alpha$ and $\beta$.

The first case is when $\alpha \neq \beta$. Rearranging the equation,
we obtain the two distinct equations.
\begin{align*}
    a_{n+1} - \alpha a_n &= \beta (a_n - \alpha a_{n-1}) \\
    a_{n+1} - \beta a_n &= \alpha (a_n - \beta a_{n-1})
\end{align*}
Because the equations from a geometric sequence, we can telescope
to obtain the following.
\begin{align*}
    a_n - \alpha a_{n-1} &= \beta^{n-1} (a_2 - \alpha a_1) \\
    a_n - \beta a_{n-1} &= \alpha^{n-1} (a_2 - \beta a_1) \\
    \therefore
    a_{n-1} = \frac{1}{\beta - \alpha} (
        \beta^{n-1} &(a_2 - \alpha a_1) -
        \alpha^{n-1} (a_2 - \beta a_1)
    )
\end{align*}
From here it is not difficult to notice that
$a_n = \frac{1}{\beta - \alpha} (\beta^n (a_2 - \alpha a_1) -
\alpha^n (a_2 - \beta a_1))$ and it won't be too difficult to
find $a_1$ and $a_2$ because they are usually given by the problem.

The second case is when $\alpha = \beta$. This case is actually
nicer because we can simply substitute the values to obtain the
following.
\begin{align*}
    a_{n+1} - 2\alpha a_n + \alpha^2 a_{n-1} &= 0 \\
    a_{n+1} - \alpha a_n &= \alpha (a_n - \alpha a_{n-1})
\end{align*}
Again, because this is a geometric sequence, we can telescope
to obtain the following equation.
\[
a_{n+1} - \alpha a_n = \alpha^{n-1} (\alpha_2 - \alpha a_1)
\]
Continuing, we can divided both sides by $\alpha^{n+1}$ to telescope
again and obtain the following.
\begin{align*}
    \frac{a_{n+1}}{\alpha^{n+1}} - \frac{a_n}{\alpha^n}
    &= \frac{a_2 - \alpha a_1}{\alpha^2} \\
    \therefore
    \frac{a_n}{\alpha^n}
    &= \frac{a_2 - \alpha a_1}{\alpha^2} \cdot (n - 1)
        + \frac{a_1}{\alpha}
\end{align*}
Therefore, $a_n = \alpha^{n-2} (n-1) (a_2 - \alpha a_1) +
\alpha^{n-1} a_1$.

This is it for this case. Honestly, I don't think its worth
memorizing the equation from this point but having the idea
to split the middle term and telescoping would be enough!
There are few more cases left, and let's continue with Case VII.

{\Large\textcolor{BrickRed}
{\textbf{VII. $\mathbf{a_n = p a_{n-1} + q b_{n-1},
b_n = r a_{n-1} + s b_{n-1}}$}}}

This is the case when we are introduced two independent variables.



\end{document}

x_{n+1}  = 6x_n + 5y_n
y_{n+1}  =  x_n + 2y_n
x_1 = 11, y_1 = 1
x_n, y_n?

